% 数据直接画为格表；每个三行矩阵的行依次对应输入的三个位置。
\begin{tikzpicture}[x=1mm,y=1mm,font=\figurefont\footnotesize,text=NNInk,
  nn flow/.append style={draw=NNWire,rounded corners=1.5mm},
  cell/.style={inner sep=0pt,minimum height=8mm},
  pe calculation/.style={align=center,inner sep=1pt,minimum height=13mm}]
  \node[nn heading] at (56,14) {输入词元序列};
  \foreach \x/\word in {38/猫,50/追,62/猫}{
    \path[fill=NNInput!18] (\x,0) rectangle ++(12,8);
    \node[cell] at ({\x+6},4) {\word};
  }
  \draw[NNInput,line width=1.1pt] (38,0) rectangle (74,8);
  \foreach \x in {50,62}{\draw[NNInput,line width=.5pt] (\x,0)--(\x,8);}

  \node[nn heading] at (26,-14) {词表编号 $u_i$};
  \node[nn heading] at (86,-14) {位置序号 $i$};
  \foreach \c/\id/\position in {0/7/1,1/12/2,2/7/3}{
    \path[fill=NNInput!12] ({8+12*\c},-29) rectangle ++(12,8);
    \path[fill=NNHidden!12] ({68+12*\c},-29) rectangle ++(12,8);
    \node[cell] at ({14+12*\c},-25) {$\id$};
    \node[cell] at ({74+12*\c},-25) {$\position$};
  }
  \draw[NNInput,line width=1.1pt] (8,-29) rectangle (44,-21);
  \draw[NNHidden,line width=1.1pt] (68,-29) rectangle (104,-21);
  \foreach \x in {20,32}{\draw[NNInput,line width=.5pt] (\x,-29)--(\x,-21);}
  \foreach \x in {80,92}{\draw[NNHidden,line width=.5pt] (\x,-29)--(\x,-21);}
  \draw[nn flow,-,sharp corners] (56,0)--(56,-25);
  \draw[nn flow] (56,-25)--(44,-25);
  \draw[nn flow] (56,-25)--(68,-25);
  \fill[NNWire] (56,-25) circle[radius=.45mm];

  \node[pe calculation] (lookup) at (26,-43) {查嵌入表 $\bm E$\\$\bm e_i^{\top}=\bm E_{u_i,:}$};
  \node[pe calculation] (encode) at (86,-43) {计算 $\bm p_i^{\top}$\\$\sin i,\ \cos i$\\$\sin(0.01i),\ \cos(0.01i)$};
  \draw[nn flow] (26,-29)--(lookup.north);
  \draw[nn flow] (86,-29)--(encode.north);
  \draw[nn flow] (lookup.south)--(26,-61);
  \draw[nn flow] (encode.south)--(86,-61);

  % 两个矩阵没有外层容器；轮廓就是矩阵的边界。
  \foreach \r in {0,1,2}{
    \pgfmathtruncatemacro{\posTint}{10+6*\r}
    \path[fill=NNInput!18] (6.8,{-69-8*\r}) rectangle ++(38.4,8);
    \path[fill=NNHidden!\posTint] (66.8,{-69-8*\r}) rectangle ++(38.4,8);
    \node[anchor=east,inner sep=0pt] at (4.8,{-65-8*\r}) {$\the\numexpr\r+1\relax$};
    \node[anchor=east,inner sep=0pt] at (64.8,{-65-8*\r}) {$\the\numexpr\r+1\relax$};
  }
  \foreach \r/\a/\b/\c/\d in {0/0.20/-0.10/0.40/0.30,1/0.00/0.50/-0.20/0.10,2/0.20/-0.10/0.40/0.30}{
    \foreach \x/\val in {11.6/\a,21.2/\b,30.8/\c,40.4/\d}{\node at (\x,{-65-8*\r}) {$\val$};}
  }
  \foreach \r/\a/\b/\c/\d in {0/0.84/0.54/0.01/1.00,1/0.91/-0.42/0.02/1.00,2/0.14/-0.99/0.03/1.00}{
    \foreach \x/\val in {71.6/\a,81.2/\b,90.8/\c,100.4/\d}{\node at (\x,{-65-8*\r}) {$\val$};}
  }
  \foreach \base/\ink in {6.8/NNInput,66.8/NNHidden}{
    \draw[\ink,line width=1.1pt] (\base,-85) rectangle ++(38.4,24);
    \foreach \c in {1,2,3}{\draw[\ink,line width=.5pt] ({\base+9.6*\c},-85)--({\base+9.6*\c},-61);}
    \foreach \y in {-69,-77}{\draw[\ink,line width=.5pt] (\base,\y)--({\base+38.4},\y);}
  }

  \node[nn pointwise] (add) at (56,-105) {$+$};
  \draw[nn flow] (26,-85)--(26,-105)--(add.west);
  \draw[nn flow] (86,-85)--(86,-105)--(add.east);
  \node[nn annotation] at (56,-97) {逐元素相加};
  \draw[nn flow] (add.south)--(56,-116);
  \foreach \r in {0,1,2}{
    \path[fill=NNOutput!18] (32,{-124-8*\r}) rectangle ++(48,8);
    \node[anchor=east,inner sep=0pt] at (29,{-120-8*\r}) {$\the\numexpr\r+1\relax$};
  }
  \foreach \r/\a/\b/\c/\d in {0/1.04/0.44/0.41/1.30,1/0.91/0.08/-0.18/1.10,2/0.34/-1.09/0.43/1.30}{
    \foreach \x/\val in {38/\a,50/\b,62/\c,74/\d}{\node at (\x,{-120-8*\r}) {$\val$};}
  }
  \draw[NNOutput,line width=1.1pt] (32,-140) rectangle (80,-116);
  \foreach \x in {44,56,68}{\draw[NNOutput,line width=.5pt] (\x,-140)--(\x,-116);}
  \foreach \y in {-124,-132}{\draw[NNOutput,line width=.5pt] (32,\y)--(80,\y);}
  \node[nn heading] at (56,-148) {输入表示 $\bm H^{(0)}\in\mathbb R^{3\times4}$};
\end{tikzpicture}
